\chapter{Introdução} \label{Cap:Introdução}

O mercado financeiro é intrinsecamente volátil, impulsionado não apenas por fundamentos econômicos, mas também por eventos políticos, macroeconômicos, regulatórios e pela percepção coletiva (sentimento) disseminada em dados não estruturados massivos, como notícias de imprensa e mídias sociais.

Embora a relação entre divulgação de resultados corporativos, indicadores macroeconômicos e movimentação de preços seja amplamente documentada, a influência do sentimento público capturado em notícias e mídias sociais sobre o comportamento dos ativos permanece menos compreendida e quantificada. O volume massivo e a alta velocidade (\textit{volume} e \textit{velocity}) de geração desses dados textuais não estruturados configuram um desafio típico de \textit{Big Data}, excedendo a capacidade de análise humana em tempo real e criando assimetrias informacionais no mercado.

Estudos empíricos demonstram que termos com conotação de incerteza nas edições diárias do jornal Valor Econômico apresentam associação negativa estatisticamente significante com o retorno diário do Ibovespa, enquanto palavras com tom negativo correlacionam-se positivamente com o aumento da volatilidade do índice \cite{galdi2018}, evidenciando o impacto quantificável do sentimento noticioso sobre a dinâmica de preços no mercado brasileiro.


\section{Problema}

Nesse contexto, emerge a seguinte questão de pesquisa: como estruturar e quantificar, em escala, a polaridade e o conteúdo semântico de publicações financeiras no X/Twitter para caracterizar o sentimento do mercado brasileiro?

Estudos internacionais já demonstraram que o sentimento expresso em mídias sociais mantém relação mensurável com o comportamento do mercado financeiro - \citeonline{bollen2011}, por exemplo, identificaram associação estatisticamente significante entre o humor coletivo expresso no Twitter e as variações do mercado de ações norte-americano. No entanto, análises robustas e validadas dessa relação para o mercado brasileiro ainda são escassas. As particularidades linguísticas do português brasileiro, a estrutura diferenciada do mercado de capitais nacional e a limitada disponibilidade de datasets rotulados de qualidade representam desafios específicos que dificultam a transferência direta de soluções desenvolvidas para outros contextos \cite{santos2023finbert}.

Iniciativas recentes, como o desenvolvimento do FinBERT-PT-BR \cite{santos2023finbert}, representam avanços importantes ao treinar modelos baseados na arquitetura BERT com 1,4 milhão de textos financeiros em português, porém ainda enfrentam limitações quanto ao tamanho das bases de dados rotuladas (apenas 503 textos anotados manualmente\footnote{A base de dados de 503 textos anotados manualmente é descrita em \citeonline{santos2023finbert}. O dataset não está disponível publicamente de forma independente; o modelo resultante pode ser acessado em: \url{https://huggingface.co/lucas-leme/FinBERT-PT-BR}.}).


\section{Justificativa}

A literatura nacional sobre análise de sentimento aplicada ao mercado financeiro permanece incipiente \cite{galdi2018}, com poucos trabalhos dedicados ao desenvolvimento de ferramentas de PLN específicas para o domínio financeiro em português brasileiro. Essa lacuna é agravada pela limitada disponibilidade de datasets rotulados de qualidade, pelas particularidades linguísticas do português e pela estrutura diferenciada do mercado de capitais nacional, fatores que dificultam a simples transposição de soluções consolidadas em outros idiomas e mercados \cite{santos2023finbert}.

Iniciativas recentes, como o FinBERT-PT-BR \cite{santos2023finbert}, representam avanços relevantes, porém ainda enfrentam restrições quanto ao tamanho das bases de dados anotadas manualmente e à ausência de pipelines reproduzíveis que integrem coleta, pré-processamento e classificação de sentimento em um fluxo estruturado. O presente trabalho busca contribuir nessa direção, ao propor e documentar um pipeline de PLN de ponta a ponta voltado ao contexto financeiro brasileiro, oferecendo à comunidade acadêmica um protocolo metodológico replicável e passível de extensão em pesquisas futuras.

Do ponto de vista prático, a capacidade de monitorar e quantificar o sentimento disseminado em publicações de veículos especializados em tempo quase real representa um recurso relevante para analistas e gestores que buscam incorporar informações não estruturadas à sua leitura de mercado. A automação desse processo por meio de técnicas de aprendizado de máquina supera a limitação humana de processar grandes volumes de dados textuais com alta velocidade, reduzindo assimetrias informacionais e ampliando a qualidade da análise qualitativa do ambiente econômico.


\section{Objetivos}\label{sec:objetivos}

\subsection{Objetivo Geral}

Desenvolver e validar pipeline de Processamento de Linguagem Natural (PLN) para análise de sentimento de publicações financeiras no X/Twitter, caracterizando padrões semânticos do discurso financeiro brasileiro.

\subsection{Objetivos Específicos}

\begin{enumerate}
	\item Coletar e estruturar um dataset de publicações do X/Twitter dos veículos InfoMoney e InvestingBrasil, abrangendo o período de outubro de 2025 a fevereiro de 2026;
	
	\item Realizar Análise Exploratória dos Dados (AED) para identificação de dados redundantes e valores atípicos (\textit{outliers}) no dataset coletado;
	
	\item Desenvolver um pipeline de pré-processamento textual e classificação de sentimento baseado em PLN para determinação da polaridade dos textos coletados;
	
	\item Avaliar e comparar o desempenho de diferentes abordagens de classificação de sentimento --- baseadas em léxico e em modelos de linguagem pré-treinados --- no domínio financeiro em português brasileiro;
	
	\item Caracterizar descritivamente o dataset resultante, identificando a distribuição de polaridade e os padrões semânticos predominantes no discurso financeiro dos veículos analisados.
	
\end{enumerate}


\section{Delimitação do Escopo}

Como \textit{proxy} para captura de sentimento de mercado, serão utilizadas publicações da rede social X (anteriormente Twitter) de veículos especializados em economia e finanças: InfoMoney e InvestingBrasil. O período de análise compreenderá publicações entre 10 de outubro de 2025 e 6 de fevereiro de 2026, totalizando aproximadamente quatro meses de dados.

A delimitação desse recorte temporal decorre de um conjunto de restrições metodológicas e operacionais: (i) o acesso à API do X/Twitter está condicionado a planos comerciais com janelas históricas limitadas, inviabilizando a coleta retroativa de grandes volumes de dados sem custos elevados; e (ii) o custo de acesso a planos de API com maior capacidade de recuperação histórica representa uma barreira concreta para pesquisas acadêmicas sem financiamento dedicado. Reconhece-se que esse recorte constitui uma limitação do estudo, cujos impactos serão discutidos no Capítulo \ref{Cap:Conclusoes}.

A escolha por posts em detrimento de artigos completos fundamenta-se em aspectos metodológicos: (i) maior facilidade de extração estruturada de dados via API, 
dispensando técnicas complexas de \textit{web scraping}; e (ii) formato textual conciso e objetivo, mais adequado para análise de polaridade semântica.


\section{Estrutura do Trabalho}

Este trabalho está organizado em cinco capítulos. O Capítulo \ref{Cap:Introdução} apresenta a contextualização do problema, a questão de pesquisa, a identificação do gap na literatura nacional, a justificativa e os objetivos geral e específicos da pesquisa.

O Capítulo \ref{Cap:Fundamentacao_Teorica} apresenta a revisão de literatura, abordando os fundamentos de Processamento de Linguagem Natural (PLN), as principais abordagens de análise de sentimento aplicadas ao mercado financeiro e os trabalhos relevantes desenvolvidos para o contexto brasileiro e para a língua portuguesa.

O Capítulo \ref{Cap:Metodologia} descreve em detalhes a metodologia proposta, incluindo: o processo de coleta de dados via API do X/Twitter; a análise exploratória e o tratamento do dataset; e o pipeline de PLN para extração de features semânticas e classificação de polaridade dos textos coletados.

O Capítulo \ref{Cap:Ava_Exp} apresenta os resultados da análise de sentimento realizada, incluindo a caracterização descritiva do dataset, os padrões semânticos identificados no discurso financeiro dos veículos analisados e a discussão dos achados à luz da literatura.

Por fim, o Capítulo \ref{Cap:Conclusoes} sintetiza as principais conclusões do trabalho, discute as contribuições científicas e práticas da pesquisa, apresenta as limitações encontradas e sugere direcionamentos para trabalhos futuros.